\subsection{Многотактные арифметические модули}
\label{sec:multi_cycle_modules}

\subsubsection{Постановка задачи}

Для базовой реализации арифметических модулей естественно было бы разделить их на три независимых модуля, каждый из которых будет реализовывать свою операцию. Однако в дальнейшем эти модули должны быть объединены в единую многотактную схему вычисления функции \eqref{eq:func} с заданными ограничениями на количество сумматоров и умножителей. Поэтому такой вариант разделения на модули не подойдёт.

Посчитаем количество базовых операционных элементов, которое требует каждая из операций согласно блок-схемам из приложения \ref{app:algorithms}:
\begin{itemize}
    \item суммирование результатов двух предыдущих операций:
    \begin{itemize}
        \item $N$-разрядный сумматор.
    \end{itemize}
    \item вычисление квадратного корня числа $b$:
    \begin{itemize}
        \item вычитатель $N$-разрядного числа;
        \item $N$-разрядное логическое ИЛИ.
    \end{itemize}
    \item вычисление кубического корня числа $a$:
    \begin{itemize}
        \item $N$-разрядное логическое ИЛИ;
        \item вычитатель $N$-разрядного числа;
        \item $N$-разрядный умножитель.
    \end{itemize}
\end{itemize}

В данной работе вычитатель будет реализован с помощью сумматора, поэтому на него распространяются ограничения на количество сумматоров.

Как можно видеть, если реализовать каждую из операций в виде отдельного модуля, то ограничения не будут соблюдены. Поэтому было принято решение вынести сумматор в отдельный общий модуль, который будет последовательно использоваться как общий ресурс.

Итого получим следующий граф зависимости модулей, который представлен на рисунке \ref{fig:module_dependencies_basic}.

\begin{figure}[h]
    \centering
    \begin{tikzpicture}[
        fsm/.style={
            draw, rectangle, thick, fill=blue!10, 
            minimum width=4cm, minimum height=1cm, 
            align=center, rounded corners
        },
        module/.style={
            draw, rectangle, thick, fill=green!10, 
            minimum width=3.5cm, minimum height=1.2cm, 
            align=center
        },
        resource/.style={
            draw, rectangle, thick, fill=orange!15, 
            minimum width=3.5cm, minimum height=1cm, 
            align=center
        },
        data/.style={
            -{Stealth[scale=1.2]}, thick, draw=black
        },
        node distance=1.5cm and 1.5cm
    ]

    % --- Узлы ---

    % Модули алгоритмов
    \node[module, xshift=-2.5cm] (cbrt) {Модуль\\кубического корня};
    \node[module, xshift=2.5cm] (sqrt) {Модуль\\квадратного корня};

    % Общие аппаратные ресурсы
    \node[resource, yshift=-2cm] (adder) {Модуль сумматора};
    \node[resource, yshift=-2cm, xshift=-5cm] (mult)  {Модуль умножителя};

    % --- Связи ---

    \draw[data] (cbrt) -- (adder);
    \draw[data] (sqrt) -- (adder);
    \draw[data] (cbrt) -- (mult);
    \end{tikzpicture}
    \caption{Граф зависимости модулей для арифметических операций.}
    \label{fig:module_dependencies_basic}
\end{figure}

\subsubsection{Модуль сумматора}

Граф переходов для модуля сумматора представлен на рисунке \ref{fig:adder_fsm}.

\begin{figure}[h]
    \caption{Граф переходов для модуля сумматора.}
    \label{fig:adder_fsm}
    \centering
\end{figure}

Код модуля сумматора представлен в листинге \ref{lst:summer} приложения \ref{app:listings}.

\subsubsection{Модуль умножителя}

Граф переходов для модуля умножителя представлен на рисунке \ref{fig:multiplier_fsm}.

\begin{figure}[h]
    \caption{Граф переходов для модуля умножителя.}
    \label{fig:multiplier_fsm}
    \centering
\end{figure}

Код модуля умножителя представлен в листинге \ref{lst:multiplier} приложения \ref{app:listings}.

\subsubsection{Модуль квадратного корня}

Граф переходов для модуля квадратного корня представлен на рисунке \ref{fig:sqrt_fsm}.

\begin{figure}[h]
    \caption{Граф переходов для модуля квадратного корня.}
    \label{fig:sqrt_fsm}
    \centering
\end{figure}

Код модуля квадратного корня представлен в листинге \ref{lst:sqrt} приложения \ref{app:listings}.

\subsubsection{Модуль кубического корня}

Граф переходов для модуля кубического корня представлен на рисунке \ref{fig:cbrt_fsm}.

\begin{figure}[h]
    \caption{Граф переходов для модуля кубического корня.}
    \label{fig:cbrt_fsm}
    \centering
\end{figure}

Код модуля кубического корня представлен в листинге \ref{lst:cbrt} приложения \ref{app:listings}.